% GENERATED FILE. Edit canonical lessons, then run npm run generate:books.
% canonical-source-hash: fnv1a64:53755d448cc1abe0
% canonical-lessons: TA-C381-tecam, TA-C381-akati, TA-C381-cinnam, TA-C381-onrinai, TA-C381-piratiniti

\chapter{One Nation}
\label{ch:ta-one-nation}
\hlchaptermodality{381}

\begin{chapteropening}
\textbf{By the end of this chapter:} \emph{I can talk about a nation, its symbols and its representatives.}

Complete the last lesson of chapter 381: I can talk about a nation, its symbols and its representatives.
\end{chapteropening}

\section[\texorpdfstring{\ta{தேசம்} (tēcam)}{தேசம் (tēcam)}]{\ta{தேசம்} (tēcam) — a nation}

\label{lesson:TA-C381-tecam}



\begin{quote}
Before the new one: say the Tamil for to contribute, then the Tamil for to cooperate.
\end{quote}

\subsection*{You'll want to know: \ta{தேசம்}}
\textbf{\ta{தேசம்}} — \emph{tēcam} — \textquotedblleft{}a nation\textquotedblright{}.

\begin{quote}
\ta{அவர்} \ta{தேசத்துக்காகத்} \ta{தன்} \ta{வாழ்க்கையைக்} \ta{கொடுத்தார்}.
\end{quote}

\begin{quote}
\emph{avar tēcattukkākat taṉ vāḻkkaiyaik koṭuttār.}
\end{quote}

\begin{quote}
He gave his life for the nation.
\end{quote}

\ta{தேசிய}, national: \ta{தேசியக்} \ta{கொடி}, the national flag.

\begin{cousinweb}
\textbf{\ta{தேசம்}} (\emph{tēcam}) is borrowed from Sanskrit \emph{deśa}, \textquotedblleft{}point, region; province, country\textquotedblright{} (CDIAL 6547). Tamil writes the ś as \ta{ச}.
\end{cousinweb}

\begin{grammarlens}[title={What you've built}]
One more word for giving reasons and opinions.
\end{grammarlens}

\subsection*{Your turn}
\begin{itemize}
\raggedright
  \item \emph{Say it:} \emph{tēcam}
  \item \emph{Say it:} \emph{tēcam}, once more
  \item \emph{Say it:} \emph{avar tēcattukkākat taṉ vāḻkkaiyaik koṭuttār.} — He gave his life for the nation.
  \item \emph{Recall:} say the Tamil for to contribute, then the Tamil for to cooperate, then say \emph{tēcam} again
\end{itemize}

\begin{tcolorbox}[breakable,colback=teal!4,colframe=teal!35!black,title={Before you move on}]
What does \ta{தேசம்} mean? (\textquotedblleft{}A nation\textquotedblright{}.) Say it once more.
\end{tcolorbox}
\section[\texorpdfstring{\ta{அகதி} (akati)}{அகதி (akati)}]{\ta{அகதி} (akati) — a refugee}

\label{lesson:TA-C381-akati}



\begin{quote}
Before the new one: say the Tamil for to cooperate, then the Tamil for a nation.
\end{quote}

\subsection*{You'll want to know: \ta{அகதி}}
\textbf{\ta{அகதி}} — \emph{akati} — \textquotedblleft{}a refugee\textquotedblright{}.

\begin{quote}
\ta{போரால்} \ta{பல} \ta{குடும்பங்கள்} \ta{அகதிகளாக} \ta{இங்கே} \ta{வந்தன}.
\end{quote}

\begin{quote}
\emph{pōrāl pala kuṭumpaṅkaḷ akatikaḷāka iṅkē vantaṉa.}
\end{quote}

\begin{quote}
Because of the war, many families came here as refugees.
\end{quote}

Plural: \ta{அகதிகள்}, \emph{akatikaḷ}.

\ta{அகதியாக}, as a refugee.

\begin{cousinweb}
\textbf{\ta{அகதி}} (\emph{akati}) is borrowed from Sanskrit \emph{agati}, \textquotedblleft{}not going; without resource, helpless\textquotedblright{} (Monier-Williams): one who has no place to go.
\end{cousinweb}

\begin{grammarlens}[title={What you've built}]
One more word for giving reasons and opinions.
\end{grammarlens}

\subsection*{Your turn}
\begin{itemize}
\raggedright
  \item \emph{Say it:} \emph{akati}
  \item \emph{Say it:} \emph{akati}, once more
  \item \emph{Say it:} \emph{pōrāl pala kuṭumpaṅkaḷ akatikaḷāka iṅkē vantaṉa.} — Because of the war, many families came here as refugees.
  \item \emph{Recall:} say the Tamil for to cooperate, then the Tamil for a nation, then say \emph{akati} again
\end{itemize}

\begin{tcolorbox}[breakable,colback=teal!4,colframe=teal!35!black,title={Before you move on}]
What does \ta{அகதி} mean? (\textquotedblleft{}A refugee\textquotedblright{}.) Say it once more.
\end{tcolorbox}
\section[\texorpdfstring{\ta{சின்னம்} (ciṉṉam)}{சின்னம் (ciṉṉam)}]{\ta{சின்னம்} (ciṉṉam) — a symbol; an emblem}

\label{lesson:TA-C381-cinnam}



\begin{quote}
Before the new one: say the Tamil for a nation, then the Tamil for a refugee.
\end{quote}

\subsection*{You'll want to know: \ta{சின்னம்}}
\textbf{\ta{சின்னம்}} — \emph{ciṉṉam} — \textquotedblleft{}a symbol; an emblem\textquotedblright{}.

\begin{quote}
\ta{புறா} \ta{அமைதியின்} \ta{சின்னம்}.
\end{quote}

\begin{quote}
\emph{puṟā amaitiyiṉ ciṉṉam.}
\end{quote}

\begin{quote}
The dove is a symbol of peace.
\end{quote}

X-\ta{இன்} \ta{சின்னம்}, a symbol of X. In elections, each party has its \ta{சின்னம்}.

\begin{cousinweb}
\textbf{\ta{சின்னம்}} (\emph{ciṉṉam}) is borrowed from Sanskrit \emph{cihna}, \textquotedblleft{}mark, sign\textquotedblright{} (CDIAL 4833); Pali has \emph{cinha}. Tamil reshapes the cluster hn as \ta{ன்ன}.
\end{cousinweb}

\begin{grammarlens}[title={What you've built}]
One more word for giving reasons and opinions.
\end{grammarlens}

\subsection*{Your turn}
\begin{itemize}
\raggedright
  \item \emph{Say it:} \emph{ciṉṉam}
  \item \emph{Say it:} \emph{ciṉṉam}, once more
  \item \emph{Say it:} \emph{puṟā amaitiyiṉ ciṉṉam.} — The dove is a symbol of peace.
  \item \emph{Recall:} say the Tamil for a nation, then the Tamil for a refugee, then say \emph{ciṉṉam} again
\end{itemize}

\begin{tcolorbox}[breakable,colback=teal!4,colframe=teal!35!black,title={Before you move on}]
What does \ta{சின்னம்} mean? (\textquotedblleft{}A symbol; an emblem\textquotedblright{}.) Say it once more.
\end{tcolorbox}
\section[\texorpdfstring{\ta{ஒன்றிணை} (oṉṟiṇai)}{ஒன்றிணை (oṉṟiṇai)}]{\ta{ஒன்றிணை} (oṉṟiṇai) — to unite; to bring together}

\label{lesson:TA-C381-onrinai}



\begin{quote}
Before the new one: say the Tamil for a refugee, then the Tamil for a symbol.
\end{quote}

\subsection*{You'll want to know: \ta{ஒன்றிணை}}
\textbf{\ta{ஒன்றிணை}} — \emph{oṉṟiṇai} — \textquotedblleft{}to unite; to bring together\textquotedblright{}.

\begin{quote}
\ta{அந்தத்} \ta{தலைவர்} \ta{மக்களை} \ta{ஒன்றிணைத்தார்}.
\end{quote}

\begin{quote}
\emph{antat talaivar makkaḷai oṉṟiṇaittār.}
\end{quote}

\begin{quote}
That leader united the people.
\end{quote}

Past (I …): \ta{ஒன்றிணைத்தேன்}, \emph{oṉṟiṇaittēṉ}.

Transitive with -\ta{ஐ} and the past in -\ta{த்த}-; people unite, \ta{ஒன்றிணைந்தார்கள்}, with -\ta{ந்த}-.

\begin{cousinweb}
\textbf{\ta{ஒன்றிணை}} (\emph{oṉṟiṇai}) is native: \emph{oṉṟu}, \textquotedblleft{}one\textquotedblright{} (DEDR 990), and \emph{iṇai}, \textquotedblleft{}to join, unite\textquotedblright{} (DEDR 457): to join into one.
\end{cousinweb}

\begin{grammarlens}[title={What you've built}]
One more word for giving reasons and opinions.
\end{grammarlens}

\subsection*{Your turn}
\begin{itemize}
\raggedright
  \item \emph{Say it:} \emph{oṉṟiṇai}
  \item \emph{Say it:} \emph{oṉṟiṇai}, once more
  \item \emph{Say it:} \emph{antat talaivar makkaḷai oṉṟiṇaittār.} — That leader united the people.
  \item \emph{Recall:} say the Tamil for a refugee, then the Tamil for a symbol, then say \emph{oṉṟiṇai} again
\end{itemize}

\begin{tcolorbox}[breakable,colback=teal!4,colframe=teal!35!black,title={Before you move on}]
What does \ta{ஒன்றிணை} mean? (\textquotedblleft{}To unite; to bring together\textquotedblright{}.) Say it once more.
\end{tcolorbox}
\section[\texorpdfstring{\ta{பிரதிநிதி} (piratiniti)}{பிரதிநிதி (piratiniti)}]{\ta{பிரதிநிதி} (piratiniti) — a representative}

\label{lesson:TA-C381-piratiniti}



\begin{quote}
Before the new one: say the Tamil for a symbol, then the Tamil for to unite.
\end{quote}

\subsection*{You'll want to know: \ta{பிரதிநிதி}}
\textbf{\ta{பிரதிநிதி}} — \emph{piratiniti} — \textquotedblleft{}a representative\textquotedblright{}.

\begin{quote}
\ta{ஒவ்வொரு} \ta{வகுப்பிலிருந்தும்} \ta{ஒரு} \ta{பிரதிநிதி} \ta{கூட்டத்துக்கு} \ta{வந்தார்}.
\end{quote}

\begin{quote}
\emph{ovvoru vakuppiliruntum oru piratiniti kūṭṭattukku vantār.}
\end{quote}

\begin{quote}
A representative from each class came to the meeting.
\end{quote}

X-\ta{இன்} \ta{பிரதிநிதியாக}, as the representative of X.

\begin{cousinweb}
\textbf{\ta{பிரதிநிதி}} (\emph{piratiniti}) is borrowed from Sanskrit \emph{pratinidhi}, \textquotedblleft{}substitution; a substitute, representative, proxy\textquotedblright{} (Monier-Williams).
\end{cousinweb}

\begin{grammarlens}[title={What you've built}]
One more word for giving reasons and opinions.
\end{grammarlens}

\subsection*{Your turn}
\begin{itemize}
\raggedright
  \item \emph{Say it:} \emph{piratiniti}
  \item \emph{Say it:} \emph{piratiniti}, once more
  \item \emph{Say it:} \emph{ovvoru vakuppiliruntum oru piratiniti kūṭṭattukku vantār.} — A representative from each class came to the meeting.
  \item \emph{Recall:} say the Tamil for a symbol, then the Tamil for to unite, then say \emph{piratiniti} again
\end{itemize}

\begin{tcolorbox}[breakable,colback=teal!4,colframe=teal!35!black,title={Before you move on}]
What does \ta{பிரதிநிதி} mean? (\textquotedblleft{}A representative\textquotedblright{}.) Say it once more.
\end{tcolorbox}
